% !PRINTED 507
\weberpage{507}

Denn alle Permutationen aus $P$, die die Ziffer $0$ ungeändert
lassen, bilden einen Theiler $Q$ von $P$. Ausserdem giebt es wegen
der Transitivität in $P$ Permutationen $\pi_1, \pi_2 \ldots \pi_{m-1}$, die die
Ziffer $0$ in $1, 2 \ldots m-1$ überführen und es ist

\[
P = Q + Q\pi_1 + Q\pi_2 + \cdots + Q\pi_{m-1},
\]

also der Grad von $P$ gleich dem $m$fachen des Grades von $Q$.

\begin{center}
§. 155.

\gesperrt{Reduction der Galois'schen Resolvente durch Ad-\\
junction. Normaltheiler einer Gruppe.}
\end{center}

Wir haben nun die allmähliche Reduction der Gruppe einer
gegebenen Gleichung zu betrachten, die durch die Adjunction
von gewissen algebraischen Grössen eintreten kann.

Es sei, wie im vorigen Paragraphen, $P$ die \gesperrt{Galois}'sche
Gruppe vom Grade $p$ und $Q$ einer ihrer Theiler vom Grade $q$
und vom Index $j$, ferner $\psi$ eine zu $Q$ gehörige Function der
Wurzeln und

\weq{1}{\psi, \psi_1, \psi_2 \ldots \psi_{j-1}}

seien die conjugirten Grössen. Wir stellen den folgenden Satz
an die Spitze.

1. \gesperrt{Die conjugirten Grössen (1) sind die Wurzeln
einer irreducibeln Gleichung vom Grade $j$
in $\Omega$.}

Denn bedeutet $t$ eine Veränderliche, so bleibt

\weq{2}{\varphi(t) = (t - \psi)(t - \psi_1) \ldots (t - \psi_{j-1})}

(nach §.154, 5.) ungeändert, wenn eine Permutation aus $P$ angewandt
 wird; also sind die Coëfficienten der Function $\varphi(t)$, deren
Wurzeln die Grössen (1) sind, in $\Omega$ enthalten [§. 149, b)].

Um die Irreducibilität von $\varphi(t)$ nachzuweisen, nehmen wir
an, es sei $\Phi(t)$ irgend eine Function in $\Omega$, die für $t = \psi$ verschwindet,
also $\Phi(\psi) = 0$. Da auf diese Gleichung alle Permutationen
 von $P$ angewandt werden dürfen, so folgt, dass auch
$\Phi(\psi_1), \Phi(\psi_2) \ldots \Phi(\psi_{j-1})$ Null sein müssen, dass also $\Phi(t)$
durch $\varphi(t)$ theilbar sein muss. Darin aber ist die Irreducibilität
enthalten.
